% 第三组统一排版模板：复制到对应报告目录后填写正文。
% 编译：latexmk -g -xelatex -outdir=_build main.tex
\documentclass[UTF8,a4paper,zihao=-4,fontset=fandol]{ctexart}
\usepackage[margin=2.5cm]{geometry}
\usepackage{amsmath,amssymb,amsfonts,amsthm}
\usepackage{graphicx,booktabs}
\usepackage{caption}
\usepackage[hidelinks]{hyperref}

\setmainfont{lmroman10-regular.otf}[
  BoldFont=lmroman10-bold.otf,
  ItalicFont=lmroman10-italic.otf,
  BoldItalicFont=lmroman10-bolditalic.otf
]
\setsansfont{lmsans10-regular.otf}[BoldFont=lmsans10-bold.otf]
\setmonofont{lmmono10-regular.otf}
\linespread{1.3}
\setlength{\parindent}{2\ccwd}
\setlength{\parskip}{0pt}
\setcounter{secnumdepth}{2}
\ctexset{
  section = {
    format = \raggedright\heiti\zihao{4},
    number = \arabic{section},
    beforeskip = 12pt,
    afterskip = 6pt
  },
  subsection = {
    format = \raggedright\heiti\zihao{-4},
    number = \thesection.\arabic{subsection},
    beforeskip = 6pt,
    afterskip = 3pt
  }
}
\pagestyle{plain}
\numberwithin{equation}{section}
\numberwithin{figure}{section}
\numberwithin{table}{section}
\DeclareCaptionFont{groupcaption}{\zihao{5}}
\captionsetup{font=groupcaption,labelsep=quad}

\theoremstyle{definition}
\newtheorem{theorem}{定理}[section]
\newtheorem{lemma}[theorem]{引理}
\newtheorem{proposition}[theorem]{命题}
\newtheorem{definition}[theorem]{定义}
\newtheorem{remark}[theorem]{注}
\newtheorem{example}[theorem]{例}
\newtheorem{conjecture}[theorem]{猜想}
\newtheorem{problem}[theorem]{问题}
\renewcommand{\proofname}{证明}

\title{\heiti\zihao{2}候选题讨论模板}
\author{\zihao{-4}第三组\\\zihao{-4}作者：待填写}
\date{}

\begin{document}
\maketitle

本文件提供组内候选题讨论的排版与填写结构。
具体题目和试解结果应在仓库转为私有后填写。

\section{题目}
\label{sec:problem}

填写题目名称和完整题面，说明要求证明、构造、计算或比较什么。
写出基础目标及完成标准，使其他成员能够据此独立开始解题。

\section{条件}
\label{sec:conditions}

定义研究对象和符号，列明假设、参数范围、边界情况、输入输出，
以及数据、计算资源和方法限制。说明题目与课程内容的联系。

\section{问题现状}
\label{sec:background}

说明背景与研究价值、已有结论、方法或基线，以及本题具体要推进的问题。
引用真实来源并注明必要的定理编号或页码；仍待查证的判断应明确标注。

\section{初步用 AI 尝试的解答结果}
\label{sec:ai}

\subsection{尝试过程与结果}

记录工具、可知模型版本、日期、实际输入的题面及关键提示，
概述主要解题思路和得到的解答或阶段性结果。
必要推导、代码和实验记录应有可访问的位置。

\subsection{人工核查与剩余问题}

写明核查人、检查步骤及结果，指出发现的错误和仍未解决的环节。
根据试解情况，说明题目的条件、目标或难度是否需要调整。

\section{后续进阶问题}
\label{sec:extensions}

至少提出一个完整、明确的进阶问题，说明相对基础题改变的条件或目标，
给出研究价值、预期难点、完成标准及可行性依据。
注明它属于本次课程的进阶目标，还是长期探索方向。

% 有真实引用后，在正文使用 \cite{引用键}，并启用以下两行。
% \bibliographystyle{unsrt}
% \bibliography{references}
% 必要附录放在参考文献后，用 \appendix 和 \section 开始。
\end{document}
